\section{Transformer Language Model Architecture}
% ======================================================================

\problem{linear}{Implementing the linear module (1 point)}

\deliverable{Implement a \texttt{Linear} class that inherits from \texttt{torch.nn.Module} and performs a linear transformation.}

\textit{Implementation summary:} 
I implemented a bias-free \texttt{Linear} module with a trainable weight matrix of shape $(d_{\text{out}}, d_{\text{in}})$. 
Its forward pass contracts the input's final dimension with the weight's input dimension to produce outputs of shape $(\ldots, d_{\text{out}})$; 
the weights are initialized with the prescribed truncated normal distribution. The adapter loads the test-provided matrix into the module before evaluating it.


\problem{embedding}{Implement the embedding module (1 point)}

\deliverable{Implement an \texttt{Embedding} class that inherits from \texttt{torch.nn.Module}.}

\textit{Implementation summary:} 
I implemented \texttt{Embedding} with a trainable matrix whose rows correspond to vocabulary entries and whose columns are embedding dimensions. 
I initialize the matrix with the prescribed truncated normal distribution. The forward pass uses token IDs to select rows, so an input of shape $(\ldots)$ 
produces vectors of shape $(\ldots, d_{\text{model}})$. The adapter loads the supplied embedding matrix before performing the lookup.


\problem{rmsnorm}{Root Mean Square Layer Normalization (1 point)}

\deliverable{Implement \texttt{RMSNorm} as a \texttt{torch.nn.Module}.}

\textit{Implementation summary:} % TODO: briefly describe your implementation


\problem{positionwise\_feedforward}{Implement the position-wise feed-forward network (2 points)}

\deliverable{Implement the SwiGLU feed-forward network, composed of a SiLU activation and a gated linear unit.}

\textit{Implementation summary:} % TODO: briefly describe your implementation


\problem{rope}{Implement RoPE (2 points)}

\deliverable{Implement a class \texttt{RotaryPositionalEmbedding} that applies RoPE to the input tensor.}

\textit{Implementation summary:} % TODO: briefly describe your implementation


\problem{softmax}{Implement softmax (1 point)}

\deliverable{Write a function to apply the softmax operation on a tensor.}

\textit{Implementation summary:} % TODO: briefly describe your implementation


\problem{scaled\_dot\_product\_attention}{Implement scaled dot-product attention (5 points)}

\deliverable{Implement the scaled dot-product attention function.}

\textit{Implementation summary:} % TODO: briefly describe your implementation


\problem{multihead\_self\_attention}{Implement causal multi-head self-attention (5 points)}

\deliverable{Implement causal multi-head self-attention as a \texttt{torch.nn.Module}.}

\textit{Implementation summary:} % TODO: briefly describe your implementation


\problem{transformer\_block}{Implement the Transformer block (3 points)}

\deliverable{Transformer block code that passes the provided tests.}

\textit{Implementation summary:} % TODO: briefly describe your implementation


\problem{transformer\_lm}{Implementing the Transformer LM (3 points)}

\deliverable{A Transformer LM module that passes the provided tests.}

\textit{Implementation summary:} % TODO: briefly describe your implementation


\problem{transformer\_accounting}{Transformer LM resource accounting (5 points)}

\textbf{(a)} How many trainable parameters would a GPT-2 XL-sized model have? How much memory is required to load it (single-precision)?

\deliverable{A one-to-two sentence response.}

\textit{Answer:} % TODO: your answer here

\textbf{(b)} Identify the matrix multiplies required for a forward pass. How many FLOPs do they require in total?

\deliverable{A list of matrix multiplies (with descriptions), and the total number of FLOPs.}

\textit{Answer:} % TODO: your answer here

\textbf{(c)} Which parts of the model require the most FLOPs?

\deliverable{A one-to-two sentence response.}

\textit{Answer:} % TODO: your answer here

\textbf{(d)} Repeat the analysis with GPT-2 small, medium, and large. How do proportional FLOPs change with model size?

\deliverable{For each model, provide a breakdown of model components and associated FLOPs (as a proportion of total). Include a one-to-two sentence description.}

\textit{Answer:} % TODO: your answer here

\textbf{(e)} Take GPT-2 XL and increase the context length to 16{,}384. How does the total FLOPs and relative contribution change?

\deliverable{A one-to-two sentence response.}

\textit{Answer:} % TODO: your answer here


